\section{\texorpdfstring{Понятие вычета аналитической функции в её
изолированной особой точке. Формула для нахождения вычета в полюсе
порядка \(m \in \mathbb{N}\). Вычет функции в точке
\(z=\infty\).}{Понятие вычета аналитической функции в её изолированной особой точке. Формула для нахождения вычета в полюсе порядка m \textbackslash in \textbackslash mathbb\{N\}. Вычет функции в точке z=\textbackslash infty.}}\label{ux43fux43eux43dux44fux442ux438ux435-ux432ux44bux447ux435ux442ux430-ux430ux43dux430ux43bux438ux442ux438ux447ux435ux441ux43aux43eux439-ux444ux443ux43dux43aux446ux438ux438-ux432-ux435ux451-ux438ux437ux43eux43bux438ux440ux43eux432ux430ux43dux43dux43eux439-ux43eux441ux43eux431ux43eux439-ux442ux43eux447ux43aux435.-ux444ux43eux440ux43cux443ux43bux430-ux434ux43bux44f-ux43dux430ux445ux43eux436ux434ux435ux43dux438ux44f-ux432ux44bux447ux435ux442ux430-ux432-ux43fux43eux43bux44eux441ux435-ux43fux43eux440ux44fux434ux43aux430-m-in-mathbbn.-ux432ux44bux447ux435ux442-ux444ux443ux43dux43aux446ux438ux438-ux432-ux442ux43eux447ux43aux435-zinfty.}

Пусть \(a\in \overline{\mathbb{C}}\) - изолированная особая точка
\(\Leftrightarrow \exists \dot{U}_{\rho}(a)\), в которой \(f(z)\)
аналитическая. Если \(a=\infty\), то
\(U_{\rho}=\{ z\in \mathbb{C}:|z|>\rho \}\) - окрестность бесконечности.

\textbf{Определение:} \emph{Вычетом} функции \(f(z)\) в точке \(a\),
обозначим \(\underset{a}{\mathrm{res}} f\), называется число \[
\frac{1}{2\pi i}\int_{\gamma_{r}^{+}(a)}f(z)\,dz=\underset{a}{\mathrm{res}} f, \ 0<r<\rho, \forall r
\] (не зависит от \(r\)).

\textbf{Замечание:} \[
\underset{a}{\mathrm{res}} f=c_{-1}
\]

\textbf{Теорема:} Формула для нахождения вычета в полюсе порядка
\(m\in \mathbb{N}\): Пусть \(a\) - полюс порядка \(m\in \mathbb{N}\).
Тогда \(\underset{a}{\mathrm{res}} f\) равен: \[
\underset{a}{\mathrm{res}}f=  \frac{1}{(m-1)!}\cdot \lim_{ z \to a }  \frac{d^{m-1}}{(dz)^{m-1}}((z-a)^{m}\cdot f(z))
\]

Вычеты в точке \(z=a=\infty\) \(\to\) \(f(z)\) аналитическая в
\(U_{\rho}(\infty)=\{ z\in \mathbb{C}:|z|>\rho \}\) - Кольцо
\(K_{\rho,\infty}(0)\implies\) ряд Лорана в точке \(a=\infty\): \[
f(z)=\sum_{n=-\infty}^{+\infty}c_{n}z^{n},
\] где
\(c_{n}= \dfrac{1}{2\pi i}\int\limits_{\gamma_{r}(0)} \dfrac{f(z)}{z^{n+1}}\,dz, \forall r>\rho\)

\textbf{Определение:} Вычетом функции \(f(z)\) в точке \(z=a=\infty\),
обозначим \(\underset{\infty}{\mathrm{res}}f\), называют число, равное:
\[
\frac{1}{2\pi i}\int_{\gamma_{r}^{-}(0)}f(z)\,dz=\underset{\infty}{\mathrm{res}} f
\]

\textbf{Замечание:} \[
\underset{\infty}{\mathrm{res}} f=-c_{-1}
\]

\textbf{Формула:} Для \(\underset{\infty}{\mathrm{res}} f\), если
\(a=\infty\) - особая устранимая точка: \[
\underset{\infty}{\mathrm{res}} f=\lim_{ z \to \infty } (z\cdot(f(\infty)-f(z)))
\]
